\documentclass{llncs}
\usepackage{graphicx}
\usepackage{listings}
\usepackage{xcolor}
\usepackage{hyperref}
\usepackage{booktabs}
\usepackage{natbib}
\bibliographystyle{plainnat}

% TODO note style, plain text, no tcolorbox, matches D07 convention
\newcommand{\TODO}[1]{\textbf{[TODO: #1]}}

\lstset{
  basicstyle=\ttfamily\small,
  columns=fullflexible,
  breaklines=true,
  frame=single,
  numbers=left,
  numberstyle=\tiny,
}

\begin{document}

\title{Design and Evaluation of a Hybrid Web Application Firewall with Lightweight Machine Learning-Assisted Detection}
\subtitle{D08: Development of Beta Version and Progress Research Paper}

\author{Arehone B. Mbadaliga (219044112)}
\institute{Academy of Computer Science and Software Engineering, University of Johannesburg}

\maketitle

\begin{abstract}
This document is the D08 Beta Version deliverable for the Hybrid Web
Application Firewall (WAF) project. It builds on the D07 Alpha Version,
which delivered a working detection pipeline evaluated offline against a
held out test set, and adds the missing piece the Alpha did not have, a
real reverse proxy that binds a port and forwards live HTTP traffic. This
document reports that a full reverse proxy module was built for the first
time in the Beta phase, tested against two live demonstration backends,
DVWA and OWASP Juice Shop, and found to correctly forward legitimate
traffic and block confirmed attacks. It also reproduces the Alpha's core
evaluation numbers from a completely clean virtual machine, confirming the
results are not tied to one developer's original setup, and reports a
verified expansion of the rule engine from 36 to 45 rules, shown to
directly catch attacks that previously relied on the machine learning
layer alone. The remaining
sections respond directly to supervisor and moderator feedback received
at the D07 presentation, covering dataset composition, the reasoning
behind every reported finding, and an honest account of where rule based
detection alone falls short of the hybrid approach.
\end{abstract}

\clearpage
\section{Introduction and Progression from D07}

D07 delivered the Alpha Version of the Hybrid WAF. It established a
six-module pipeline (M1 TrafficReader through M6 Logger), trained three
machine learning pipelines on a merged CSIC 2010 and HTTPParams dataset
of 128{,}132 requests, and evaluated four configurations offline against
a held out test set of 25{,}627 requests: rule-only, ML-only, hybrid, and
ModSecurity with the OWASP Core Rule Set v4 as an external baseline. The
hybrid configuration reached an F1 score of 0.826, against 0.449 for the
rule engine alone and 0.568 for ModSecurity.

What D07 did not have was a live reverse proxy. Every D07 result came
from feeding a CSV of pre-extracted requests through the pipeline objects
directly in Python, and the interactive demonstration used for the D07
presentation, \texttt{demo/interactive\_demo.py}, called the same pipeline
objects directly and never opened a network socket. This is stated plainly
here because it matters for what counts as new work in this document:
the reverse proxy described in Section~\ref{sec:m1} did not exist before
the Beta phase.

D08 has four aims, matching the deliverable's stated focus and the
feedback raised by the supervisor and moderator at the D07 presentation:

\begin{enumerate}
  \item Build and test a real reverse proxy, so the WAF can sit in front
        of an actual web application rather than only being evaluated
        offline (Section~\ref{sec:m1}).
  \item Confirm the Alpha's results reproduce from a clean environment,
        not just the machine they were first measured on
        (Section~\ref{sec:reproduce}).
  \item Respond directly to the D07 feedback: explain how each finding
        was produced, characterise the dataset and what drives its
        attack composition, and defend the gap between rule-only and
        hybrid detection with concrete examples
        (Sections~\ref{sec:findings}--\ref{sec:defense}).
  \item Expand detection coverage in the rule engine and feature
        extractor, and record honestly which of that expansion is
        complete and which is still in progress
        (Section~\ref{sec:expansion}).
\end{enumerate}

\clearpage
\section{Response to D07 Feedback}
\label{sec:feedback}

Table~\ref{tab:feedback} maps each point raised at the D07 presentation
to the section of this document that answers it, following the same
practice used in D04 and D06 of tracking feedback explicitly rather than
folding it in unlabelled.

\begin{table}[h]
\centering
\caption{D07 presentation feedback and corresponding responses in this document}
\label{tab:feedback}
\begin{tabular}{p{0.45\textwidth}p{0.45\textwidth}}
\toprule
\textbf{Feedback point} & \textbf{Response in this document} \\
\midrule
Explain all findings and how they were found &
Section~\ref{sec:findings} traces every reported number back to the
exact script that produced it \\
Consider more keywords on the feature extractor &
Section~\ref{sec:expansion} records the current F3 keyword list and the
expansion still in progress \\
What contributes to the attacks based on the datasets? &
Section~\ref{sec:dataset} breaks down the merged dataset by source and
category \\
Explain the journey from catching traffic to decision &
Section~\ref{sec:m1} walks a single request through all six modules \\
Dataset characteristics, distribution, imbalance, training and results &
Section~\ref{sec:dataset} and Section~\ref{sec:reproduce} \\
Defend rule-based misses versus the hybrid difference &
Section~\ref{sec:defense} \\
Shallow SQL keywords, consider IP tagging &
Section~\ref{sec:expansion} records this as planned, not yet built \\
Explain param={payload} clearly &
Section~\ref{sec:param} \\
\bottomrule
\end{tabular}
\end{table}

\clearpage
\section{The Journey: Traffic to Decision}
\label{sec:m1}

\subsection{Why this section exists}
The moderator asked for a clear account of how a request moves through
the system from the moment it arrives to the moment a decision is made.
D07 could not fully answer this, since nothing in the Alpha actually
received live traffic. This section describes the module built for
Beta that makes this a real, testable journey rather than a description
of six components that had never been wired together end to end.

\subsection{M1 TrafficReader: built new for Beta}
\texttt{src/traffic\_reader.py} did not exist before this deliverable.
It is a \texttt{BaseHTTPRequestHandler} subclass, following the interface
named in the original D04 proposal (Section 4.2), that binds
\texttt{127.0.0.1:8080} and threads each incoming connection so multiple
requests can be handled concurrently.

\subsection{The path of one request}
Take a single request, a GET to \texttt{/vulnerabilities/sqli/} carrying
the query string \texttt{id=1' OR '1'='1\&Submit=Submit}, sent as a
normal browser would send it, URL-encoded.

\begin{enumerate}
  \item \textbf{M1 TrafficReader} accepts the connection and reads the
        method, path and query string. The query string arrives
        URL-encoded (\texttt{1\%27\%20OR\%20\%271\%27\%3D\%271}) and is
        decoded once, here, before any other module sees it.
        \TODO{This decoding step was added after a real bug was found
        during Beta testing: without it, the rule engine could not match
        an encoded quote character. Confirm this stays flagged as a
        genuine finding in the final write-up, not smoothed over.}
  \item \textbf{M2 FeatureExtractor} computes the six-feature vector
        (F1 length, F2 special character count, F3 SQL keyword count,
        F4 script flag, F5 traversal count, F6 Shannon entropy) from the
        decoded text.
  \item \textbf{M3 RuleEngine} checks the decoded text against all 36
        loaded regex rules. For this request, rule SQLi\_002 matches.
  \item \textbf{M4 MLClassifier} scores the same feature vector through
        the trained Random Forest, independently of the rule engine's
        result.
  \item \textbf{M5 DecisionEngine} combines the two verdicts under the
        active operating mode. In hybrid mode, if both layers agree the
        request is an attack, the decision is reported as \emph{Both
        agree: confirmed attack}. If only the classifier flags it, the
        decision is reported as Case B, an obfuscated payload the rule
        layer missed.
  \item \textbf{M6 Logger} records the request, the feature vector, both
        verdicts, the final decision and the per-stage timings to
        \texttt{logs/waf\_events.csv}.
  \item If the decision is ALLOW, M1 opens a connection to the
        configured backend, forwards the original method, path, query
        string, headers and body, and relays the backend's response back
        to the client unchanged. If the decision is BLOCK, M1 returns an
        HTTP 403 with the decision's stated reason, and the backend is
        never contacted.
\end{enumerate}

\subsection{Verification}
This flow was tested against two real backends running as separate
Docker containers, following the same "intercepted backend" placement
described in D06: DVWA on port 8081 and OWASP Juice Shop on port 3000,
never both targeted at once.

\begin{itemize}
  \item A clean request to DVWA's \texttt{/index.php} was forwarded
        correctly, returning DVWA's real 302 redirect to its login page
        with session cookies intact.
  \item The SQLi payload above was blocked with the reason \emph{Both
        agree: confirmed attack}, and DVWA never received the request.
  \item A clean request to Juice Shop's \texttt{/rest/products/search}
        endpoint was forwarded correctly, returning Juice Shop's real
        JSON response, confirming the backend target is genuinely
        configurable and not hardcoded to one application.
\end{itemize}

\clearpage
\section{Deployment Environment for Beta}
\label{sec:env}

The guest VM is confirmed as Ubuntu Server 24.04.4 LTS (kernel
6.8.0-124-generic). This settles an inconsistency between two earlier
documents: D04 stated Ubuntu 22.04 LTS, D06's deployment diagram stated
24.04 LTS. The 24.04 figure is correct, confirmed directly with
\texttt{lsb\_release -a} on the running guest.

\begin{table}[h]
\centering
\caption{Confirmed software versions on the Beta deployment VM}
\begin{tabular}{ll}
\toprule
\textbf{Component} & \textbf{Version} \\
\midrule
Guest OS & Ubuntu Server 24.04.4 LTS \\
Python & 3.12.3 \\
Docker & docker.io 29.1.3-0ubuntu3~24.04.2 \\
numpy & 2.5.3 \\
scikit-learn & 1.9.1 (models trained under 1.9.0, see note below) \\
scipy & 1.18.1 \\
PyYAML & 6.0.3 \\
pytest & 9.1.1 \\
\bottomrule
\end{tabular}
\end{table}

One honest note belongs here. The trained model files were pickled
under scikit-learn 1.9.0, but \texttt{requirements.txt} specifies only
\texttt{scikit-learn>=1.3.0} with no upper bound, so a fresh install on
the Beta VM pulled 1.9.1. This produced an \texttt{InconsistentVersion\allowbreak
Warning} at load time. The reproduced results in
Section~\ref{sec:reproduce} were not affected in any way that changed
the reported metrics, but this is a real reproducibility risk that
should not recur, so \texttt{requirements.txt} is being pinned to the
exact training version as part of Beta.

DVWA and Juice Shop run as Docker containers on the guest VM, matching
the ports committed to in D06's deployment diagram, 8081 and 3000
respectively. Both are now set with a \texttt{restart: unless-stopped}
policy so they come back up automatically after the VM reboots, since
manual restarts were found to be an easy step to forget between working
sessions.

\clearpage
\section{Reproducibility of Alpha Results}
\label{sec:reproduce}

To answer the concern that D07's results might be specific to the
original development machine, the full evaluation pipeline was rerun on
a completely fresh clone of the repository on the Beta VM: a fresh
\texttt{git clone}, a fresh Python virtual environment, and the raw
CSIC 2010 and HTTPParams source files copied in unmodified. No code was
changed for this test.

\texttt{build\_dataset.py} reproduced the exact same dataset composition
reported in D07:

\begin{table}[h]
\centering
\caption{Dataset build, reproduced on a clean VM}
\begin{tabular}{lrr}
\toprule
\textbf{Source} & \textbf{Requests} & \textbf{Label} \\
\midrule
CSIC 2010, normalTrafficTraining.txt & 36{,}000 & benign \\
CSIC 2010, normalTrafficTest.txt & 36{,}000 & benign \\
CSIC 2010, anomalousTrafficTest.txt & 25{,}065 & attack \\
HTTPParams & 31{,}067 & mixed \\
\midrule
\textbf{Combined total} & \textbf{128{,}132} & 91{,}304 benign, 36{,}828 attack \\
\bottomrule
\end{tabular}
\end{table}

\texttt{evaluate\_configurations.py} was then rerun against the same
held out test split (25{,}627 requests) using the already trained model
files committed to the repository, no retraining. Table~\ref{tab:reproduce}
compares the original D07 numbers against this Beta reproduction.

\begin{table}[h]
\centering
\caption{D07 original results versus Beta reproduction on a clean VM}
\label{tab:reproduce}
\begin{tabular}{lrrrr}
\toprule
\textbf{Configuration} & \textbf{Precision} & \textbf{Recall} & \textbf{F1} & \textbf{FPR} \\
\midrule
Rule-only (D07) & 1.000 & 0.289 & 0.449 & 0.000 \\
Rule-only (Beta reproduction) & 1.000 & 0.2893 & 0.4488 & 0.0000 \\
\addlinespace
ML-only, Random Forest (D07) & 0.980 & 0.713 & 0.825 & 0.006 \\
ML-only, Random Forest (Beta reproduction) & 0.9800 & 0.7126 & 0.8252 & 0.0059 \\
\addlinespace
Hybrid (D07) & 0.980 & 0.714 & 0.826 & 0.006 \\
Hybrid (Beta reproduction) & 0.9800 & 0.7135 & 0.8258 & 0.0059 \\
\addlinespace
ModSecurity + CRS v4 (D07) & 0.799 & 0.440 & 0.568 & 0.045 \\
ModSecurity + CRS v4 (Beta reproduction) & 0.7991 & 0.4401 & 0.5676 & 0.0446 \\
\bottomrule
\end{tabular}
\end{table}

Every figure reproduces to within rounding of its D07 value. This is
worth stating plainly rather than glossed over: the small differences in
the third or fourth decimal place come from floating point summation
order and not from any change in code, data, or trained model, and they
do not alter any conclusion drawn in D07.

\clearpage
\section{Dataset Characteristics and Attack Contribution}
\label{sec:dataset}

The merged dataset combines two public sources rather than one, for a
reason worth stating directly, a single source risks the model learning
that source's particular formatting habits rather than the underlying
attack patterns. CSIC 2010 contributes 97{,}065 requests
(36{,}000 + 36{,}000 benign, 25{,}065 attack), drawn from a
Spanish e-commerce site's traffic and its associated anomalous traffic
set. HTTPParams contributes 31{,}067 requests, drawn from
parameter-level injection samples rather than full HTTP requests.

The combined class balance is 91{,}304 benign to 36{,}828 attack, a
ratio of roughly 2.48 to 1. This imbalance is preserved through the
stratified 80/20 train-test split (\texttt{random\_state=42}), so the
held out test set of 25{,}627 requests keeps the same proportion. This
matters directly for how the reported recall figures should be read:
recall of 0.714 for the hybrid configuration means just over seven in
ten of the roughly 7{,}366 attack requests in the test set were caught,
not seven in ten of all traffic.

\subsection{What contributes to the attacks}

The moderator asked directly what drives the attack composition of the
dataset, how much is SQLi versus XSS versus path traversal versus
something else. Answering this honestly requires two different methods,
kept separate here rather than blended together, because they are not
the same kind of evidence.

\textbf{HTTPParams} carries a real \texttt{attack\_type} column in its
source CSV, genuine ground truth assigned when the dataset was built,
not inferred by this project. Table~\ref{tab:httpparams} reports it
directly.

\begin{table}[h]
\centering
\caption{HTTPParams attack composition, ground truth from the dataset's own \texttt{attack\_type} column}
\label{tab:httpparams}
\begin{tabular}{lrr}
\toprule
\textbf{Category} & \textbf{Requests} & \textbf{Share of attacks} \\
\midrule
SQLi & 10{,}852 & 92.3\% \\
XSS & 532 & 4.5\% \\
Path traversal & 290 & 2.5\% \\
Command injection & 89 & 0.8\% \\
\midrule
\textbf{Total} & \textbf{11{,}763} & 100.0\% \\
\bottomrule
\end{tabular}
\end{table}

\textbf{CSIC 2010} carries no such column. It was built as a binary
normal/anomalous dataset only, so there is no ground truth subtype to
report for its 25{,}065 attack rows. Rather than guess, each attack
row's own text was run through the real 45-rule RuleEngine (the same
engine used in production, see Section~\ref{sec:expansion}), and
categorised by whichever rule fired. Rows that matched no rule at all
are reported as uncategorised, since that is itself real, meaningful
information about this dataset, not a gap to hide.
Table~\ref{tab:csic} reports this.

\begin{table}[h]
\centering
\caption{CSIC 2010 attack composition, inferred via the RuleEngine (no ground truth subtype exists in this dataset)}
\label{tab:csic}
\begin{tabular}{lrr}
\toprule
\textbf{Category} & \textbf{Requests} & \textbf{Share of attacks} \\
\midrule
Uncategorised (no rule matched) & 24{,}125 & 96.2\% \\
SQLi & 754 & 3.0\% \\
XSS & 186 & 0.7\% \\
Path traversal & 0 & 0.0\% \\
\midrule
\textbf{Total} & \textbf{25{,}065} & 100.0\% \\
\bottomrule
\end{tabular}
\end{table}

This result matters beyond simply answering the moderator's question.
96.2\% of CSIC 2010's attack corpus does not match any of the 45 rules
in the rule engine, including the 9 added during Beta. This is not a
weakness in the rule engine's coverage in the ordinary sense, CSIC
2010's anomalous traffic was deliberately constructed to include attacks
that do not follow simple textbook syntax, exactly the kind of traffic
the rule layer was never going to catch by design (see the structural
argument in Section~\ref{sec:defense}). It is direct, concrete evidence
for why this project uses a hybrid approach at all: the majority of
CSIC 2010's attack traffic can only be caught by a layer that reasons
statistically over the whole request rather than matching a fixed
pattern, exactly the role the Random Forest classifier plays through
Case B.

\clearpage
\section{How Each Finding Was Produced}
\label{sec:findings}

This section exists because the moderator asked, directly, for every
reported number to be traceable to the script that produced it, not
just stated as a result. Table~\ref{tab:provenance} lists that mapping.

\begin{table}[h]
\centering
\caption{Provenance of every reported metric}
\label{tab:provenance}
\begin{tabular}{p{0.32\textwidth}p{0.32\textwidth}p{0.26\textwidth}}
\toprule
\textbf{Result} & \textbf{Script} & \textbf{What it measures} \\
\midrule
Dataset composition (Section~\ref{sec:dataset}) & \texttt{build\_dataset.py} &
Parses raw source files, merges, extracts features, writes
\texttt{data/merged\_dataset.csv} \\
Four-configuration comparison (Sections~\ref{sec:reproduce}, \ref{sec:defense}) &
\texttt{evaluate\_configurations.py} & Scores the held out 25{,}627-row
test split through rule-only, ML-only and hybrid \\
ModSecurity baseline & \texttt{export\_test\_set\_for\_modsec.py} then
\texttt{benchmark\_modsecurity.py} & Exports the same test rows in a
format ModSecurity can be sent, then measures its real block/allow
decisions \\
Latency & \texttt{benchmark\_latency.py} & Times M1 entry to M5 exit
with \texttt{time.perf\_counter()} across the full test set \\
Unit test count & \texttt{pytest tests/ -v} & Runs the 23 unit tests
covering the FeatureExtractor, RuleEngine and DecisionEngine directly \\
\bottomrule
\end{tabular}
\end{table}

Two things are deliberately excluded from this table because they did
not produce a reported number: \texttt{demo/interactive\_demo.py}, which
is a terminal demonstration tool that prints results but is not itself
an evaluation script, and \texttt{src/traffic\_reader.py}, whose
correctness for Beta was confirmed through the manual curl tests
described in Section~\ref{sec:m1} rather than an automated benchmark.
\TODO{Consider whether an automated integration test for M1 belongs in
the plan for the final deliverable, since manual curl testing is fine
for development but weaker evidence for a paper than a repeatable
script.}

\clearpage
\section{Defending the Rule-Only Gap}
\label{sec:defense}

The rule-only configuration reaches perfect precision (1.000) but recall
of only 0.289. This means 71.1\% of the actual attacks in the test set
passed straight through undetected when only the 36 regex rules were
used. This number needs defending honestly rather than presented as a
surprise, since it is the entire justification for adding a machine
learning layer at all.

The reason is structural, not a bug. Regex rules can only match text
patterns they were written to expect. \texttt{rules.txt} currently
contains 36 rules across three categories (SQLi, XSS, PATH), each
targeting a specific known syntax, for example rule SQLi\_002 matching
the literal characters of \texttt{' OR '1'='1}. Any attack that does not
contain one of these 36 specific patterns, whether because it is
obfuscated, encoded, or simply a variant the rule author did not
anticipate, is invisible to this layer by design. \citep{durmuskaya2025}
make the same point about pattern-based detection generally, that
static rule sets cannot generalise beyond the patterns they were given.

The hybrid configuration raises recall from 0.289 to 0.714 without
changing the rule engine at all. The entire gain comes from
\emph{Case B}, the condition where the rule engine reports no match but
the trained classifier's probability score crosses the 0.5 threshold
anyway. In the reproduced Beta test run, Case B fired on 3{,}232 of
25{,}627 test requests, 12.61\%. These are, by definition, exactly the
requests the rule engine structurally cannot catch, obfuscated or
unusual attack payloads that still carry the statistical fingerprint a
Random Forest trained on the feature vector (length, special character
density, keyword count, script flag, traversal count, entropy) can
pick up even without matching a known syntax.

The rule engine was expanded during Beta from 36 to 45 rules, described
in detail in Section~\ref{sec:expansion}, and that expansion closed a
small, verifiable part of this gap, recall rose from 0.2893 to 0.2927.
The gap remains structural even after expansion though, since 45 rules
is still 45 fixed patterns, and the remaining 3{,}207 Case B detections
in the expanded configuration are, by the same reasoning, attacks no
fixed pattern was written to expect.

This is also why the hybrid configuration's false positive rate rises
slightly above rule-only's 0.000 to 0.006. A classifier making
probabilistic judgements on unseen text will occasionally flag a benign
request that happens to share statistical features with attack traffic,
where a rule engine that only matches exact known patterns essentially
cannot false-positive on anything it was not written to match. This
trade-off, a small increase in false positives for a large increase in
recall, is stated here as a deliberate, measured trade-off rather than
a weakness left unexplained.

\clearpage
\section{The param={payload} Correction, Explained}
\label{sec:param}

This finding was first reported in D07 but, based on presentation
feedback, did not land clearly enough to be understood on its own. This
section restates it with a direct before-and-after comparison.

The HTTPParams dataset provides raw payload strings, not full query
strings. An early version of \texttt{build\_dataset.py} wrapped every
HTTPParams payload as \texttt{param=\{payload\}} before feature
extraction, under the assumption this would make the payload look like
a realistic query parameter.

\begin{table}[h]
\centering
\caption{The param= wrapping bug, before and after}
\begin{tabular}{p{0.2\textwidth}p{0.35\textwidth}p{0.35\textwidth}}
\toprule
& \textbf{Before (with wrapper)} & \textbf{After (fixed)} \\
\midrule
Raw HTTPParams row & \texttt{hello world} & \texttt{hello world} \\
Text sent to FeatureExtractor & \texttt{param=hello world} & \texttt{hello world} \\
F2 special character count & 1 (the injected \texttt{=} sign) & 0 \\
Effect & Every benign row, including plain English text with no
special characters at all, was artificially given a non-zero F2
score purely because of the wrapper, not because of anything in the
original text & F2 correctly reflects the actual content of the row \\
\bottomrule
\end{tabular}
\end{table}

This mattered because F2 special character count is one of six inputs
to the classifier, and an artificial non-zero F2 on thousands of benign
rows pushed the model toward over-predicting attacks on ordinary text.
Removing the wrapper and using the raw payload directly as the query
string improved the Random Forest's precision from 0.965 to 0.980 and
its false positive rate from 1.1\% to 0.6\%. The finding is kept in this
document, as it was in D07, as a deliberate example of a real bug found
through testing and corrected honestly, not smoothed over.

\clearpage
\section{Rule Engine and Feature Extractor Expansion}
\label{sec:expansion}

\TODO{The rule engine subsection below is now complete, added, tested,
and re-evaluated. The feature extractor and IP tagging subsections that
follow are still open work, recorded here with a clear current state
rather than claimed as done, following the project's standing rule that
a design decision is only claimed once it is verified.}

\subsection{Rule engine (M3), expanded and verified}
\texttt{rules.txt} grew from 36 to 45 rules during Beta, an even split
across SQLi, XSS and PATH (15 each). The 9 additions cover stacked
query injection (\texttt{SQLi\_013}), SQL Server command execution via
\texttt{xp\_cmdshell} (\texttt{SQLi\_014}), blind boolean-based SQLi
functions such as \texttt{SUBSTRING} (\texttt{SQLi\_015}),
\texttt{data:} URI and \texttt{vbscript:} based XSS
(\texttt{XSS\_013}, \texttt{XSS\_014}), iframe \texttt{srcdoc}
injection (\texttt{XSS\_015}), and PHP stream wrapper, SSH key path,
and IIS configuration file path traversal targets
(\texttt{PATH\_013}--\texttt{PATH\_015}). Each rule has a matching unit
test (UT14 through UT22, following the same fixture pattern as the
existing UT09--UT13), and the full suite of 32 unit tests passes.

\texttt{evaluate\_configurations.py} was rerun against the same held
out 25{,}627-row test split with the expanded 45-rule set, no changes
to either trained model. Table~\ref{tab:ruleexpansion} compares the
36-rule and 45-rule results directly.

\begin{table}[h]
\centering
\caption{Effect of the 9 added rules on rule-only and hybrid detection}
\label{tab:ruleexpansion}
\begin{tabular}{lrrrr}
\toprule
\textbf{Configuration} & \textbf{TP} & \textbf{Recall} & \textbf{F1} & \textbf{Case B rate} \\
\midrule
Rule-only, 36 rules & 2{,}131 & 0.2893 & 0.4488 & 12.61\% (3{,}232) \\
Rule-only, 45 rules & 2{,}156 & 0.2927 & 0.4528 & 12.51\% (3{,}207) \\
Hybrid, 36 rules & 5{,}256 & 0.7135 & 0.8258 & --- \\
Hybrid, 45 rules & 5{,}256 & 0.7135 & 0.8258 & --- \\
\bottomrule
\end{tabular}
\end{table}

The result is worth stating plainly rather than left to the table alone.
Rule-only true positives rose by exactly 25, and Case B's count fell by
exactly the same 25 requests, from 3{,}232 to 3{,}207. This is not a
coincidence, it means the 9 new rules are directly catching the same
25 attacks that previously only the ML classifier caught. Hybrid's
overall recall and F1 did not change, because those 25 attacks were
already being caught by the hybrid system either way, through Case B.
What changed is which layer catches them, a small but real and
verifiable shift of detection work from the probabilistic ML layer to
the deterministic rule layer, which is a meaningful result in its own
right: it shows the rule expansion generalises correctly rather than
just adding rules that happen to fire on already-caught traffic or,
worse, on benign traffic (false positives on both configurations remain
at 0).

\subsection{Feature extractor (M2), current state}
\texttt{FeatureExtractor}'s F3 feature currently checks 13 SQL keywords
in the code, a wider list than the 9 originally documented in D07's
Table 5, itself an inconsistency worth correcting regardless of any
further expansion. The moderator's request for a wider keyword set is
still open. Any change to this list changes the F3 value for every one
of the 128{,}132 rows in the merged dataset, meaning both trained models
would need retraining and \texttt{evaluate\_configurations.py} rerunning
before an updated result can be reported honestly. This is why the
rule engine expansion above was prioritised first for Beta, since it
does not carry the same retraining cost.

\subsection{IP tagging, planned}
The supervisor's suggestion of tagging or tracking source IP addresses
when the SQL keyword count is low but a request still looks suspicious
has not been built. \texttt{HTTPRequest} already carries a
\texttt{source\_ip} field, so the data needed is available, but no
decision logic using it exists yet. This is recorded here as planned
work for the remainder of Beta or for the final deliverable, not as a
completed feature.

\clearpage
\section{Limitations and Planned Work}
\label{sec:limitations}

Stated plainly, the following items are known and open at the time of
this document:

\begin{itemize}
  \item The rule and feature extractor expansion described in
        Section~\ref{sec:expansion} is drafted but not implemented,
        tested, or re-evaluated.
  \item The full attack-category breakdown of the merged dataset
        (Section~\ref{sec:dataset}) has not been produced.
  \item The monitoring dashboard (M6, port 8082) has not yet been
        tested against traffic generated by the live reverse proxy,
        only against the offline evaluation logs.
  \item \texttt{requirements.txt} does not yet pin an exact
        scikit-learn version, which is a real reproducibility risk
        noted in Section~\ref{sec:env}.
  \item \texttt{datasets/README.md} still lists WADE-2022 as a dataset
        to download, which contradicts this project's own finding, kept
        from D07, that WADE-2022 could not be verified as a real,
        citable dataset. \citep{calvo2022} is the reference this
        project has been able to verify for the adaptive WAF concept
        WADE-2022 was intended to represent; the dataset itself was
        excluded and this should be corrected in the repository.
\end{itemize}

\clearpage
\section{Conclusion}

This document reports the Beta phase's central achievement, a real
reverse proxy that did not exist at D07, built, tested and confirmed
working against two live demonstration backends. It also confirms, by
independent reproduction on a clean virtual machine, that the Alpha's
offline evaluation results were not specific to one developer's original
setup. The remaining work is scoped honestly in
Section~\ref{sec:limitations} rather than presented as finished, and
directly answers the supervisor and moderator's D07 feedback as mapped
in Table~\ref{tab:feedback}.

\bibliography{references}

\end{document}
